\lysh{36654}{4}

\begin{alist}
\item Αν η εταιρεία δεν κατασκευάσει μπλουζάκια (κατασκευάσει $0$ μπλουζάκια), τότε έχουμε
$K(0) = \text{120 \euro}$. Επομένως η εταιρεία, χωρίς να κατασκευάζει μπλουζάκια, έχει (πάγια) έξοδα \text{120 \euro}.

\item Η μοναδιαία μεταβολή του $x$, δηλαδή αν τα μπλουζάκια που κατασκευάζονται αυξηθούν κατά $1$ (από $x$ γίνουν $x + 1$), θα προκαλέσει σταθερή μεταβολή \text{12,5 \euro} στο κόστος κατασκευής και σταθερή μεταβολή στα έσοδα \text{15,5 \euro}. Δηλαδή,
$K(x + 1) - K(x) = 12,5$ και $E(x + 1) - E(x) = 15,5$

\item Τα μπλουζάκια που πρέπει να πωληθούν ώστε τα έσοδα να είναι ίσα με τα έξοδα είναι η λύση της εξίσωσης $E(x) = K(x)$. Είναι:
\[E(x) = K(x) \Leftrightarrow 15,5x = 12,5x + 120 \Leftrightarrow 3x = 120 \Leftrightarrow x = 40\]
Επομένως, για να μην <<μπαίνει μέσα>> η επιχείρηση, πρέπει να πουλήσει $40$ μπλουζάκια.

\item Αν πουλήσουν $60$ μπλουζάκια τα έσοδα θα είναι
$E(60) = 15,5 \cdot 60 = 930\text{ \euro}$
και τα έξοδα $K(60) = 12,5 \cdot 60 + 120 = 870\text{ \euro}$.
Επομένως το κέρδος τους θα είναι $930 - 870 = 60\text{ \euro}$.
\end{alist}

